%% ch10.tex --- Evaluation, Bias and Variance
%% Numbers are computed by verify/ch10_bias_variance.py.
%% Dataset T10 is the same synthetic set used in the companion lecture slides.

\chapter{Evaluation, Bias and Variance}
\label{ch:ch10}

\epigraph{The first principle is that you must not fool yourself --- and you
are the easiest person to fool.}{---\ Richard Feynman}

\begin{chapterlearning}
After working through this chapter you will be able to:
\begin{enumerate}
  \item Explain why a model must be judged on data it was not fitted to, and why
        a test set used to make a decision is no longer a test set.
  \item Use a training/validation/test split to choose among models, and say
        what each part is for.
  \item Carry out $k$-fold cross-validation and say when it is worth its cost.
  \item Diagnose high bias and high variance from training and validation error,
        measured against a baseline.
  \item Derive the bias--variance decomposition of expected squared error, and
        estimate its terms by simulation.
  \item Read the effect of the regularization parameter $\lambda$ on bias and
        variance.
  \item Draw and interpret learning curves, and use them to decide whether more
        data will help.
  \item Apply the bias/variance recipe for neural networks.
  \item Check whether a classifier's probabilities are calibrated, and report the
        Brier score and log loss alongside accuracy.
\end{enumerate}
\end{chapterlearning}

Every model in the book can now be built and trained. This part of the book is
about what happens next. The first model is rarely good enough, and there are a
dozen plausible ways to improve it: more data, fewer features, more features,
polynomial features, more or less regularization, a bigger network, a longer
training run. Some of these take an afternoon and some take months, and choosing
by instinct is how teams lose a quarter.

The alternative is to \emph{measure} what is wrong first. This chapter builds the
measurements; \cref{ch:ch11} uses them to drive the day-to-day loop of model
development.

\section{The Diagnostic Mindset}
\label{sec:diagnostic}

Suppose a regularized linear model makes unacceptably large errors. Candidate
remedies include:

\begin{enumerate}
  \item collect more training examples;
  \item try a smaller set of features;
  \item collect additional features;
  \item add polynomial features;
  \item decrease $\lambda$;
  \item increase $\lambda$.
\end{enumerate}

Each is sensible for some models and useless for others; \cref{sec:what-next}
shows that they split cleanly into remedies for two different diseases.

\begin{defbox}{Diagnostic}
A \term{diagnostic} is a test run on a learning algorithm to find out what is
and is not working, so as to decide how to improve it. Diagnostics cost time to
run, but far less than pursuing the wrong remedy.
\end{defbox}

\section{Evaluating a Model}
\label{sec:evaluating}

\Cref{ch:ch07} established that training cost measures fit, not
generalization. The basic remedy is to hold data back.

\begin{defbox}{Test error}
Split the data into a training set and a \term{test set}, often $70$--$80\%$ and
$20$--$30\%$. Fit on the training set only, then evaluate
\[
  J_{\text{test}} = \frac{1}{2m_{\text{test}}}\sum_{i=1}^{m_{\text{test}}}
     \bigl(f(\xii{i}_{\text{test}}) - \yii{i}_{\text{test}}\bigr)^2
\]
\emph{without} any regularization term: the penalty is part of how the model is
fitted, not of how it is judged. For classification, the usual test measure is
the \term{misclassification error}, the fraction of test examples with
$\yhat \ne y$.
\end{defbox}

Throughout the chapter we use one synthetic dataset, whose truth we know.

\begin{defbox}{Dataset T10}
Fifty points with $x$ uniform on $[0, 3]$ and
\[
  y = \sin 2x + 0.5x + \varepsilon, \qquad \varepsilon \sim \mathcal{N}(0,\ 0.25^2),
\]
generated by
\begin{lstlisting}[style=pythonml]
rng = np.random.default_rng(4)
x = rng.uniform(0, 3, 50)
y = np.sin(2 * x) + 0.5 * x + rng.normal(0, 0.25, 50)
\end{lstlisting}
and split in order: the first $20$ for training, the next $15$ for validation,
the last $15$ for testing. The models are polynomials in $x$, fitted by least
squares on standardised powers. Because the noise has standard deviation
$0.25$, no model can do better on average than $J = 0.25^2/2 = 0.03125$.
\end{defbox}

\begin{pitfall}{A test set used twice is not a test set}
The test error is an honest estimate of performance on new data only if nothing
about the model was chosen by looking at it. The moment you compare two models
on the test set and keep the better one, the test set has become part of the
fitting process, and its error is now an optimistic estimate --- the best of
several tries. Every subsequent decision informed by it makes it worse.
\end{pitfall}

\section{Model Selection}
\label{sec:model-selection}

Choosing the polynomial degree is a decision of exactly the kind the pitfall
warns about. The fix is a third set.

\begin{defbox}{Training, validation and test sets}
Split the data three ways, for example $60/20/20$:
\begin{itemize}
  \item the \term{training set} fits the parameters $\vw$ and $b$;
  \item the \term{validation set} (also called the cross-validation set or dev
        set) chooses everything else --- the degree, $\lambda$, the
        architecture, the learning rate;
  \item the \term{test set} is used once, at the end, to report the chosen
        model's performance.
\end{itemize}
Each set has one job, and any decision taken by looking at a set turns that
set's error into an underestimate.
\end{defbox}

\begin{workedexample}{Choosing the degree on T10}
Fit polynomials of degree $1$ to $10$ to the training set and evaluate each on
the validation and test sets:

\begin{center}
\small
\renewcommand{\arraystretch}{1.15}
\begin{tabular}{@{}rrrr@{}}
\toprule
degree & $J_{\text{train}}$ & $J_{\text{cv}}$ & $J_{\text{test}}$\\
\midrule
$1$  & $0.0806$ & $0.1119$ & $0.1285$\\
$2$  & $0.0777$ & $0.0991$ & $0.1158$\\
$3$  & $0.0219$ & $\mathbf{0.0496}$ & $0.0429$\\
$4$  & $0.0209$ & $0.0515$ & $0.0444$\\
$5$  & $0.0203$ & $0.0510$ & $0.0410$\\
$6$  & $0.0203$ & $0.0516$ & $0.0408$\\
$7$  & $0.0201$ & $0.0518$ & $0.0380$\\
$8$  & $0.0167$ & $0.1472$ & $0.1121$\\
$9$  & $0.0165$ & $0.3432$ & $0.2200$\\
$10$ & $0.0164$ & $0.6085$ & $0.3362$\\
\bottomrule
\end{tabular}
\end{center}

The validation error is smallest at degree $3$, so we choose degree $3$ and
report its test error, $0.0429$, once. The training error only falls with the
degree, as \cref{sec:nested} proved it must; the validation error falls,
flattens, and rises --- the U of \cref{fig:degree-curve}.

Had we chosen on the test column instead, we would have picked degree $7$ and
reported $0.0380$. That number is not wrong arithmetic; it is the luckiest of
ten draws, from a degree the validation set rates no better than degrees $3$ to
$6$. The difference between $0.0429$ and $0.0380$ is the size of the
self-deception.
\end{workedexample}

\begin{figure}[htbp]
  \centering
  \begin{tikzpicture}
    \begin{axis}[mlaxis, width=0.8\linewidth, height=5.2cm,
      xlabel={degree $d$}, ylabel={cost}, xtick={1,...,10},
      xmin=0.5, xmax=10.5, ymin=0, ymax=0.65,
      legend style={font=\scriptsize, at={(0.03,0.97)}, anchor=north west}]
      \addplot[MLBlue, very thick, mark=*, mark size=1.4pt] coordinates
        {(1,0.0806)(2,0.0777)(3,0.0219)(4,0.0209)(5,0.0203)(6,0.0203)(7,0.0201)
         (8,0.0167)(9,0.0165)(10,0.0164)};
      \addlegendentry{$J_{\text{train}}$}
      \addplot[MLRose, very thick, mark=square*, mark size=1.4pt] coordinates
        {(1,0.1119)(2,0.0991)(3,0.0496)(4,0.0515)(5,0.0510)(6,0.0516)(7,0.0518)
         (8,0.1472)(9,0.3432)(10,0.6085)};
      \addlegendentry{$J_{\text{cv}}$}
      \addplot[MLGray, dashed, domain=0.5:10.5, samples=2] {0.03125};
      \addlegendentry{noise floor $0.03125$}
    \end{axis}
  \end{tikzpicture}
  \caption{Training and validation cost against polynomial degree on T10. On
  the left both are high: high bias. On the right the training cost is low and
  the validation cost climbs steeply: high variance. The minimum of the
  validation curve, at degree $3$, is the model to choose.}
  \label{fig:degree-curve}
\end{figure}

\section{\texorpdfstring{$k$}{k}-Fold Cross-Validation}
\label{sec:kfold}

A fixed validation set has two weaknesses. It withholds data from training ---
$15$ of T10's $35$ non-test points --- and the decision rests on one particular
split, which may be lucky or unlucky. When data is scarce, both matter.

\begin{algorithmbox}{$k$-fold cross-validation}
\begin{lstlisting}[style=pseudocode]
shuffle the non-test data; split into k equal folds
for each candidate model M:
    for i = 1 to k:
        fit M on all folds except fold i
        e[i] <- error of M on fold i
    CV(M) <- mean of e[1..k]     # report the spread too
choose the M with the smallest CV(M)
refit on all non-test data; report test error once
\end{lstlisting}
Every example is used for validation exactly once and for training $k-1$ times.
The cost is $k$ fits per candidate instead of one. Common choices are $k = 5$
and $k = 10$; $k = m$ is \emph{leave-one-out} cross-validation.
\end{algorithmbox}

\begin{workedexample}{Five folds on T10}
Pool T10's $35$ training and validation points, shuffle them, and split them
into five folds of seven. For each degree, fit five times and average:

\begin{center}
\small
\renewcommand{\arraystretch}{1.15}
\begin{tabular}{@{}rrr@{\qquad}rrr@{}}
\toprule
degree & mean & s.d. & degree & mean & s.d.\\
\midrule
$1$ & $0.1076$ & $0.0128$ & $6$  & $0.0438$ & $0.0184$\\
$2$ & $0.0977$ & $0.0162$ & $7$  & $0.0779$ & $0.0607$\\
$3$ & $\mathbf{0.0372}$ & $0.0169$ & $8$  & $0.0839$ & $0.0583$\\
$4$ & $0.0375$ & $0.0174$ & $9$  & $0.1410$ & $0.1661$\\
$5$ & $0.0383$ & $0.0201$ & $10$ & $0.1301$ & $0.1164$\\
\bottomrule
\end{tabular}
\end{center}

Cross-validation also chooses degree $3$. Refitted on all $35$ points, the
degree-$3$ model has test error $0.0390$ --- better than the $0.0429$ of the
model fitted on $20$, because it was trained on more data. The spread column is
the extra information a single split cannot give: degrees $3$, $4$ and $5$ are
within a fraction of one standard deviation of each other, so the data does not
really distinguish them, while from degree $7$ upward the error is not only
larger but erratic from fold to fold --- a symptom of high variance.
\end{workedexample}

\begin{researchspotlight}{Stone (1974) --- Cross-validatory choice and assessment}
\textbf{The problem.} How can one choose a statistical predictor and estimate
its accuracy from the same data, without the estimate being biased by the
choice?

\textbf{The idea.} Leave out part of the data, fit to the rest, and predict the
part left out; repeat over all such parts, and choose the method whose held-out
predictions are best. Stone formalised this ``cross-validatory'' procedure and
analysed it as a method for both choosing and assessing.

\textbf{The result.} A general, assumption-light framework for model selection
that needs no theory of the model's complexity --- only the ability to refit it.
The paper also warns, with characteristic clarity, that the assessment of a
procedure that includes cross-validatory choice requires a second, outer layer of
held-out data.

\textbf{Why it is here.} It is \cref{sec:kfold}. The warning is the reason the
test set stays untouched: when cross-validation chooses the degree, the
cross-validation error of the chosen degree is itself optimistic.
\parencite{stone1974}
\end{researchspotlight}

\section{Bias and Variance}
\label{sec:bias-variance}

\Cref{ch:ch07} introduced \emph{bias} and \emph{variance} as names for
underfitting and overfitting. The two errors give a working diagnosis.

\begin{center}
\small
\renewcommand{\arraystretch}{1.3}
\begin{tabular}{@{}L{2.6cm}L{2.0cm}L{3.0cm}L{3.2cm}@{}}
\toprule
diagnosis & $J_{\text{train}}$ & $J_{\text{cv}}$ & T10 example\\
\midrule
high bias (underfit) & high & high, $\approx J_{\text{train}}$ & degree $1$: $0.0806$, $0.1119$\\
high variance (overfit) & low & $\gg J_{\text{train}}$ & degree $10$: $0.0164$, $0.6085$\\
both & high & $\gg J_{\text{train}}$ & rare in one dimension; common in networks\\
neither & low & low, near $J_{\text{train}}$ & degree $3$: $0.0219$, $0.0496$\\
\bottomrule
\end{tabular}
\end{center}

``High'' and ``low'' need a reference point; \cref{sec:baseline} supplies it.
First, the theory behind the names.

\subsection{The decomposition}

Imagine repeating the whole experiment many times: draw a fresh training set,
fit the model, and predict at a fixed input $x$. The prediction $\hat f(x)$ is
then a random variable, varying with the training set.

\begin{derivation}{Bias--variance decomposition of squared error}
Let $y = f(x) + \varepsilon$ with $\Exp[\varepsilon] = 0$ and
$\operatorname{Var}(\varepsilon) = \sigma^2$, the noise independent of the
training set. Write $\bar f(x) = \Exp[\hat f(x)]$ for the average prediction
over training sets. The expected squared error at $x$, over both the training
set and the new noise, is
\[
  \Exp\bigl[(y - \hat f(x))^2\bigr]
  = \Exp\bigl[(f(x) - \hat f(x) + \varepsilon)^2\bigr]
  = \Exp\bigl[(f(x) - \hat f(x))^2\bigr] + \sigma^2 ,
\]
since the cross term $2\Exp[\varepsilon]\,\Exp[f(x) - \hat f(x)]$ vanishes by
independence and $\Exp[\varepsilon] = 0$. Now add and subtract $\bar f(x)$
inside the remaining square:
\[
  \Exp\bigl[(f - \hat f)^2\bigr]
  = \Exp\bigl[(f - \bar f + \bar f - \hat f)^2\bigr]
  = (f - \bar f)^2 + \Exp\bigl[(\bar f - \hat f)^2\bigr]
    + 2(f - \bar f)\,\Exp[\bar f - \hat f] .
\]
The last term is zero because $\Exp[\hat f] = \bar f$ by definition. Hence
\[
  \boxed{\;\Exp\bigl[(y - \hat f(x))^2\bigr]
  = \underbrace{\bigl(f(x) - \bar f(x)\bigr)^2}_{\text{bias}^2}
  + \underbrace{\Exp\bigl[(\hat f(x) - \bar f(x))^2\bigr]}_{\text{variance}}
  + \underbrace{\sigma^2}_{\text{noise}}\;}
\]
Bias is how far the \emph{average} model is from the truth; variance is how much
individual models scatter around their average; noise is what no model can
remove.
\end{derivation}

\begin{workedexample}{Measuring the three terms}
Because we know T10's generating curve, we can estimate the three terms
directly. Draw $2000$ training sets of $30$ points each from the same
distribution, fit each degree to each, and evaluate on a grid of $51$ points in
$[0.25,\ 2.75]$. Averaged over the grid (squared error without the $\frac12$,
so the noise term is $\sigma^2 = 0.0625$):

\begin{center}
\small
\renewcommand{\arraystretch}{1.15}
\begin{tabular}{@{}rrrrr@{}}
\toprule
degree & bias$^2$ & variance & sum with $\sigma^2$ & simulated error\\
\midrule
$1$ & $0.1266$ & $0.0157$ & $0.2047$ & $0.2054$\\
$2$ & $0.1119$ & $0.0258$ & $0.2002$ & $0.2004$\\
$3$ & $0.0026$ & $0.0081$ & $\mathbf{0.0732}$ & $0.0733$\\
$5$ & $0.0000$ & $0.0137$ & $0.0762$ & $0.0763$\\
$7$ & $0.0000$ & $0.0402$ & $0.1027$ & $0.1031$\\
$9$ & $0.0002$ & $0.3415$ & $0.4042$ & $0.4059$\\
\bottomrule
\end{tabular}
\end{center}

The decomposition checks: the sum of the three terms matches a direct simulation
of the expected error to the third decimal in every row. Bias falls to
essentially zero by degree $5$ --- a quintic can follow $\sin 2x + 0.5x$ on this
interval --- and variance grows rapidly beyond it. The total is smallest at
degree $3$, which is also what the validation set chose.

One detail repays attention: the variance at degree $3$ ($0.0081$) is
\emph{lower} than at degree $1$ ($0.0157$). A straight line forced onto a curve
swings considerably depending on where the training points happen to fall. The
textbook picture, with variance rising monotonically in complexity, is a
tendency, not a law.
\end{workedexample}

\begin{researchspotlight}{Geman, Bienenstock \& Doursat (1992) --- The bias/variance dilemma}
\textbf{The problem.} Neural networks and other non-parametric methods can, in
principle, learn any function given enough data. Why, then, do they perform so
poorly on many real problems?

\textbf{The idea.} Decompose the expected error of an estimator into bias and
variance, as in \cref{sec:bias-variance}, and examine both as functions of model
flexibility and sample size.

\textbf{The result.} Flexible estimators have low bias and high variance, and the
variance falls only slowly with more data; reducing it substantially requires
building in prior knowledge --- constraints, invariances, architecture --- that
deliberately introduces bias. ``The bias/variance dilemma'' became the standard
name for the trade-off.

\textbf{Why it is here.} It is the decomposition of this section, argued for
networks specifically. Its conclusion is the justification for regularization
(\cref{ch:ch07}) and for the weight sharing of convolutional layers
(\cref{ch:ch09}), both of which buy lower variance with a little bias.
\parencite{geman1992}
\end{researchspotlight}

\section{Regularization and Bias/Variance}
\label{sec:reg-bv}

The regularization parameter moves a model along the same axis as the degree,
in the opposite direction. Fit the degree-$10$ polynomial to T10 with the L2
penalty of \cref{sec:regularized-cost}:

\begin{center}
\small
\renewcommand{\arraystretch}{1.15}
\begin{tabular}{@{}rrrr@{}}
\toprule
$\lambda$ & $J_{\text{train}}$ & $J_{\text{cv}}$ & $J_{\text{test}}$\\
\midrule
$0$         & $0.0164$ & $0.6085$ & $0.3362$\\
$10^{-4}$   & $0.0174$ & $0.0749$ & $0.0602$\\
$10^{-3}$   & $0.0189$ & $0.0555$ & $0.0415$\\
$10^{-2}$   & $0.0198$ & $\mathbf{0.0538}$ & $0.0397$\\
$10^{-1}$   & $0.0211$ & $0.0560$ & $0.0461$\\
$1$         & $0.0299$ & $0.0687$ & $0.0707$\\
$10$        & $0.0525$ & $0.0965$ & $0.1083$\\
$10^{2}$    & $0.0768$ & $0.1171$ & $0.1376$\\
\bottomrule
\end{tabular}
\end{center}

Small $\lambda$ leaves the degree-$10$ model nearly free: high variance. Large
$\lambda$ forces its weights towards zero: high bias. The validation error is
smallest at $\lambda = 10^{-2}$, and the chosen model's test error is $0.0397$
--- as good as the best unregularized degree, from a model that was given ten
features and told only to keep its weights small.

\begin{keyinsight}{Two mirror-image curves}
Plotted against the degree, variance grows to the right. Plotted against
$\lambda$, variance grows to the \emph{left}. In both cases the validation curve
is U-shaped, and its minimum is chosen the same way: try a range, keep the
smallest validation error, report the test error once.
\end{keyinsight}

\section{Baselines}
\label{sec:baseline}

A training error of $10\%$ is excellent on some problems and dismal on others. A
bias/variance diagnosis needs a reference.

\begin{defbox}{Baseline level of performance}
A \term{baseline} is the error you can reasonably hope to reach: human-level
performance on the task, the error of a competing system, or the error achieved
on similar problems before. With a baseline in hand, two gaps diagnose a model:
\[
  \text{bias gap} = J_{\text{train}} - \text{baseline}, \qquad
  \text{variance gap} = J_{\text{cv}} - J_{\text{train}} .
\]
\end{defbox}

\begin{workedexample}{Speech recognition (illustrative numbers)}
Some audio is so noisy that people cannot transcribe it either. Suppose human
transcribers have a word error rate of $10.6\%$, and three systems give:

\begin{center}
\small
\renewcommand{\arraystretch}{1.15}
\begin{tabular}{@{}lrrr@{}}
\toprule
 & system A & system B & system C\\
\midrule
human error      & $10.6\%$ & $10.6\%$ & $10.6\%$\\
training error   & $10.8\%$ & $15.0\%$ & $15.0\%$\\
validation error & $14.8\%$ & $15.5\%$ & $19.7\%$\\
\midrule
bias gap         & $0.2$ & $4.4$ & $4.4$\\
variance gap     & $4.0$ & $0.5$ & $4.7$\\
\midrule
diagnosis        & variance & bias & both\\
\bottomrule
\end{tabular}
\end{center}

Without the baseline, system A's $10.8\%$ training error would look like high
bias. Against $10.6\%$, it is nearly as good as people, and the problem is
entirely in the gap to validation.
\end{workedexample}

T10 has a perfect baseline: the noise floor $0.03125$. The degree-$3$ model's
training error, $0.0219$, is \emph{below} it --- possible only because the
model has absorbed some of the training noise --- and its validation error,
$0.0496$, is above it. The remaining gap is variance, which is why more data,
rather than a more flexible model, is the remedy for T10.

\section{Learning Curves}
\label{sec:learning-curves}

\begin{defbox}{Learning curve}
A \term{learning curve} plots training and validation error against the number
of training examples, $m_{\text{train}}$, for a fixed model.
\end{defbox}

As $m_{\text{train}}$ grows, training error rises --- it is harder to fit more
points --- and validation error falls, because the model sees more of the
population. What happens next depends on the diagnosis.

\begin{figure}[htbp]
  \centering
  \begin{minipage}[t]{0.48\linewidth}
    \centering
    \begin{tikzpicture}
      \begin{axis}[mlaxis, width=\linewidth, height=4.8cm, xmode=log,
        xlabel={$m_{\text{train}}$}, ylabel={cost}, title={degree 1: high bias},
        xmin=8, xmax=500, ymin=0, ymax=0.2,
        xtick={10,20,50,100,200,400}, xticklabels={10,20,50,100,200,400},
        scaled y ticks=false, yticklabel style={/pgf/number format/fixed},
        legend style={font=\tiny, at={(0.97,0.5)}, anchor=north east}]
        \addplot[MLBlue, very thick, mark=*, mark size=1.1pt] coordinates
          {(10,0.0805)(12,0.0954)(15,0.1001)(20,0.1033)(30,0.1156)(50,0.1111)
           (100,0.1141)(200,0.1176)(400,0.1186)};
        \addlegendentry{$J_{\text{train}}$}
        \addplot[MLRose, very thick, mark=square*, mark size=1.1pt] coordinates
          {(10,0.1331)(12,0.1362)(15,0.1268)(20,0.1234)(30,0.1201)(50,0.1179)
           (100,0.1169)(200,0.1162)(400,0.1156)};
        \addlegendentry{$J_{\text{cv}}$}
        \addplot[MLGray, dashed, domain=8:500, samples=2] {0.03125};
        \addlegendentry{noise floor}
      \end{axis}
    \end{tikzpicture}
  \end{minipage}\hfill
  \begin{minipage}[t]{0.48\linewidth}
    \centering
    \begin{tikzpicture}
      \begin{axis}[mlaxis, width=\linewidth, height=4.8cm, xmode=log,
        xlabel={$m_{\text{train}}$}, title={degree 8: high variance},
        xmin=15, xmax=500, ymin=0, ymax=0.2,
        xtick={20,50,100,200,400}, xticklabels={20,50,100,200,400},
        scaled y ticks=false, yticklabel style={/pgf/number format/fixed}]
        \addplot[MLBlue, very thick, mark=*, mark size=1.1pt] coordinates
          {(20,0.0151)(30,0.0205)(50,0.0234)(100,0.0280)(200,0.0305)(400,0.0311)};
        \addplot[MLRose, very thick, mark=square*, mark size=1.1pt] coordinates
          {(20,0.1831)(30,0.0645)(50,0.0470)(100,0.0402)(200,0.0374)(400,0.0365)};
        \addplot[MLGray, dashed, domain=15:500, samples=2] {0.03125};
      \end{axis}
    \end{tikzpicture}
  \end{minipage}
  \caption{Learning curves on data from T10's distribution: the median over
  $100$ random training sets of each size, evaluated on a fixed validation set
  of $200$ points. Left: the straight line's two curves meet at about $0.117$,
  far above the noise floor, and stay there. Right: the degree-$8$ model's gap
  is large for small training sets and closes as data grows, approaching the
  floor.}
  \label{fig:learning-curves}
\end{figure}

\begin{keyinsight}{What a learning curve says about more data}
\textbf{High bias.} The curves flatten and meet, well above the baseline. From
$100$ to $400$ examples the straight line's validation error improves by
$0.0013$. Collecting more data will not help; the model cannot represent the
pattern, however much of it it sees.

\textbf{High variance.} There is a large gap that narrows as $m_{\text{train}}$
grows: the degree-$8$ validation error falls from $0.1831$ at $20$ examples to
$0.0365$ at $400$. More data helps, and the curve shows roughly how much.

A learning curve costs a handful of fits on subsets of data you already have,
and can save the expense of collecting data that would not have helped.
\end{keyinsight}

\section{Deciding What to Try Next}
\label{sec:what-next}

The six remedies of \cref{sec:diagnostic} now sort themselves:

\begin{center}
\small
\renewcommand{\arraystretch}{1.25}
\begin{tabular}{@{}L{5.0cm}L{4.0cm}@{}}
\toprule
action & treats\\
\midrule
collect more training examples & high variance\\
try a smaller set of features & high variance\\
increase $\lambda$ & high variance\\
\midrule
collect additional features & high bias\\
add polynomial features & high bias\\
decrease $\lambda$ & high bias\\
\bottomrule
\end{tabular}
\end{center}

\begin{pitfall}{Shrinking the training set to ``fix'' training error}
Fewer training examples make the training error look better --- that is the left
end of every learning curve --- but they never fix bias, and they make variance
worse. A model that does badly on its training data needs more capacity or
better features, not less data.
\end{pitfall}

\section{Bias and Variance in Neural Networks}
\label{sec:nn-bv}

For classical models, bias and variance pull against each other: a simpler model
has less variance and more bias, and the practitioner balances them. Large
neural networks loosen that trade-off, because a large enough network trained on
a modest dataset is usually a \emph{low-bias} model. That gives a recipe.

\begin{algorithmbox}{The bias/variance recipe for neural networks}
\begin{lstlisting}[style=pseudocode]
repeat:
    if training error is high (compared with the baseline):
        use a bigger network (more units, more layers)
    elif validation error >> training error:
        get more data, or regularize more
    else:
        done
\end{lstlisting}
\end{algorithmbox}

Two facts make the recipe work. A larger network, appropriately regularized,
almost never does worse than a smaller one; its cost is computation, not
accuracy. And data and computation, unlike the model's inherent capacity, can
usually be bought. The limits of the recipe are therefore practical: a network
too large to train in the time available, or data too expensive to collect ---
which is when \cref{ch:ch11}'s techniques for getting more from less come in.
In Keras, the L2 penalty of \cref{ch:ch07} is attached to a layer as
\code{kernel\_regularizer=tf.keras.regularizers.L2(0.01)}.

\begin{researchspotlight}{Domingos (2012) --- A few useful things to know about machine learning}
\textbf{The problem.} Much of what makes machine learning succeed in practice is
``folk knowledge'' that textbooks leave out, and newcomers learn it slowly and
expensively.

\textbf{The idea.} Collect twelve such lessons in one short article: that
learning is representation plus evaluation plus optimization; that
generalization is what counts; that data alone is not enough; that overfitting
has many faces; that intuition fails in high dimensions; that more data beats a
cleverer algorithm; and others.

\textbf{The result.} One of the most widely read papers in the field, valued
because each lesson is stated plainly and backed by an argument rather than
anecdote.

\textbf{Why it is here.} Its section on bias and variance --- with the warning
that a learner can overfit even when the true model is in its hypothesis space,
and underfit when strong assumptions are wrong --- is the content of this
chapter in four pages. Read it at the end of the course; it will mean more then.
\parencite{domingos2012}
\end{researchspotlight}

\section{Are the Probabilities Probabilities?}
\label{sec:calibration}

A classifier that outputs probabilities makes a claim beyond its decisions:
among the cases it scores $0.8$, about $80\%$ should be positive.
\Cref{ch:ch06} deferred the check of that claim to here.

\begin{defbox}{Calibration}
A classifier is \term{calibrated} if, among all examples it assigns probability
about $p$, a fraction about $p$ are positive. A \term{reliability table} checks
this: group the predictions into bins of predicted probability, and compare the
mean prediction in each bin with the observed fraction of positives.
\end{defbox}

\begin{workedexample}{Same accuracy, different honesty}
Draw $20{,}000$ examples with $x$ uniform on $[-3,3]$ and
$\Prob(y = 1\mid x) = \gsig(1.5x)$. Classifier A outputs exactly
$\gsig(1.5x)$, the true probability. Classifier B outputs $\gsig(4.5x)$: the
same decision boundary at $x = 0$, so the same decisions and the same accuracy,
$0.8508$, but logits three times too large.

\begin{center}
\small
\renewcommand{\arraystretch}{1.12}
\begin{tabular}{@{}lrrrr@{}}
\toprule
 & \multicolumn{2}{c}{A: calibrated} & \multicolumn{2}{c}{B: overconfident}\\
\cmidrule(lr){2-3}\cmidrule(lr){4-5}
bin & predicted & observed & predicted & observed\\
\midrule
$[0.0, 0.1)$ & $0.041$ & $0.039$ & $0.010$ & $0.100$\\
$[0.1, 0.2)$ & $0.145$ & $0.134$ & $0.147$ & $0.365$\\
$[0.4, 0.5)$ & $0.450$ & $0.443$ & $0.448$ & $0.445$\\
$[0.8, 0.9)$ & $0.853$ & $0.859$ & $0.855$ & $0.616$\\
$[0.9, 1.0]$ & $0.959$ & $0.959$ & $0.991$ & $0.898$\\
\midrule
accuracy   & \multicolumn{2}{c}{$0.8508$} & \multicolumn{2}{c}{$0.8508$}\\
Brier score & \multicolumn{2}{c}{$0.1067$} & \multicolumn{2}{c}{$0.1227$}\\
log loss    & \multicolumn{2}{c}{$0.3440$} & \multicolumn{2}{c}{$0.5587$}\\
\bottomrule
\end{tabular}
\end{center}

A's predictions agree with reality in every bin. B says $0.99$ when the truth is
$0.90$, and $0.01$ when it is $0.10$: a tenfold understatement of the minority
risk in a bin holding $8{,}432$ of the $20{,}000$ examples, invisible to
accuracy.
\end{workedexample}

Two numbers summarise what accuracy misses. The \term{Brier score} is the mean
squared difference between predicted probability and outcome,
$\frac1m\sum_i(f^{(i)} - \yii{i})^2$. The \term{log loss} is the mean logistic
loss of \cref{ch:ch06}. Both reward calibrated, confident predictions and punish
overconfident ones, and both rank A above B although the two make identical
decisions. Overconfidence is common in practice: high-capacity models trained to
near-zero training loss, such as large networks, are often overconfident, and the
separable-data pathology of \cref{ch:ch06} is its extreme form.

\begin{pitfall}{Accuracy alone for a model whose probabilities will be used}
If a model's outputs feed a decision with costs --- a threshold chosen as in
\cref{sec:threshold}, a price, a triage queue --- then the probabilities, not
just the decisions, must be right. Report a reliability table and the log loss
or Brier score alongside accuracy, and recalibrate on validation data (for
instance by fitting a one-parameter logistic regression to the logits) if the
table shows a systematic distortion.
\end{pitfall}

\begin{labbox}{Lab 10 --- Model evaluation and bias/variance}
Generate T10 with the listing of \cref{sec:evaluating} and reproduce the degree
table of \cref{sec:model-selection}: the validation error must be smallest at
degree $3$ with $J_{\text{cv}} = 0.0496$. Implement $5$-fold cross-validation
with the fold assignment of your choice and confirm that it also chooses a low
degree. Reproduce the $\lambda$ table for degree $10$.

Then draw the two learning curves of \cref{fig:learning-curves} and state, for
each, whether more data would help. Finally simulate the bias--variance
decomposition for degrees $1$, $3$ and $9$, and confirm that the three terms sum
to the directly simulated error.
\end{labbox}

\section*{Chapter Summary}
\addcontentsline{toc}{section}{Chapter Summary}

\begin{itemize}
  \item Diagnose before acting: the remedies for a poor model divide into those
        for high bias and those for high variance, and only measurement says
        which applies.
  \item Judge models on held-out data, with no regularization term in the
        evaluation. A test set used to make a choice becomes optimistic.
  \item Training, validation and test sets fit, choose and report. On T10 the
        validation set chooses degree $3$ ($J_{\text{test}} = 0.0429$); choosing
        on the test set would have reported $0.0380$ for degree $7$.
  \item $k$-fold cross-validation uses every non-test example for validation once,
        at $k$ times the cost, and reports a spread. Five folds also choose degree
        $3$ on T10, with test error $0.0390$ after refitting on all $35$ points.
  \item Expected squared error decomposes into bias$^2$, variance and noise. On
        T10's distribution the sum matches simulation to three decimals; bias
        vanishes by degree $5$ and variance reaches $0.3415$ at degree $9$.
  \item $\lambda$ trades variance for bias in the mirror image of the degree:
        $\lambda = 10^{-2}$ gives the degree-$10$ model test error $0.0397$.
  \item A baseline turns errors into a diagnosis: bias gap is training error
        minus baseline, variance gap is validation minus training error.
  \item Learning curves meet above the baseline under high bias, so more data
        will not help; under high variance the gap closes as data grows.
  \item Large networks are usually low-bias: grow the network until training
        error is acceptable, then add data or regularization until validation
        error is.
  \item Calibrated probabilities match observed frequencies. Two classifiers can
        share an accuracy of $0.8508$ and differ in log loss by $0.21$.
\end{itemize}

\section*{Exercises}
\addcontentsline{toc}{section}{Exercises}

\exercisesection{Short Questions}

\sq List six actions you might take to improve a model, and say why a diagnostic
should come first.

\sq Why must the regularization term be omitted when computing the test error?

\sq State the role of each part of a training/validation/test split.

\sq Explain why selecting the degree with the test set makes the reported test
error optimistic.

\sq What does $k$-fold cross-validation provide that a single validation split
does not, and what does it cost?

\sq Give the pattern of training and validation error that identifies high bias,
and the pattern for high variance.

\sq Define a baseline level of performance and give three ways of obtaining one.

\sq Human error is $2.0\%$, training error $2.3\%$, validation error $7.1\%$.
Compute both gaps and give the diagnosis.

\sq Sketch learning curves for a high-bias and a high-variance model, and say for
each whether more data will help.

\sq State the bias/variance recipe for neural networks.

\sq What does it mean for a classifier to be calibrated? Name two scores that
reward calibration.

\exercisesection{Analytical Problems}

\ap Derive the bias--variance decomposition of $\Exp[(y - \hat f(x))^2]$,
stating every assumption used, and identify the term that no model can reduce.

\ap For the estimator $\hat f(x) = c$ (a constant fitted by least squares as the
mean of $m$ training targets) with $y = f(x) + \varepsilon$, compute the bias and
variance at a point $x$, and show that the variance is $\sigma^2/m$ plus a term
due to the spread of $f$ over the training inputs.

\ap Ridge regression with one centred feature has
$w = \sum_i\xii{i}\yii{i}/(\sum_i(\xii{i})^2 + \lambda)$. Assuming
$y = \beta x + \varepsilon$ with fixed inputs, compute the bias and variance of
$w$ as functions of $\lambda$, and show that the bias grows and the variance
shrinks as $\lambda$ increases.

\ap Explain why a degree-$3$ polynomial can have lower variance than a straight
line on T10's distribution, as the simulation found, and construct a simpler
example in which an underfitting model has larger variance than a correctly
specified one.

\ap With $k$ folds of a dataset of size $m$, each fit uses $m(k-1)/k$ examples.
Explain the effect of $k$ on the bias of the cross-validation estimate and on its
cost, and argue why $k = 5$ or $10$ is a common compromise.

\ap Show that the Brier score of a classifier that always predicts the base rate
$\pi$ is $\pi(1-\pi)$, and use this to interpret a Brier score of $0.1067$ when
$\pi \approx 0.5$.

\exercisesection{Design Problems}

\dpr A hospital has $3000$ labelled chest X-rays from $900$ patients and wants to
choose among five architectures and three values of $\lambda$. Design the split
--- including how you keep one patient's images together --- the selection
procedure, and what you report.

\dpr A crop-yield model has training $R^2 = 0.62$ and validation $R^2 = 0.60$,
while an expert's estimate achieves $0.75$. Diagnose the model and design the
next three experiments in priority order, stating the result that would make you
stop each one.

\dpr Collecting $1000$ more labelled examples will cost a large sum. Design an
experiment, using only the data you already have, that predicts whether it is
worth paying for.

\dpr A network takes four hours to train and you have one week of compute.
Design a procedure to choose $\lambda$ and the network width, justify your grid,
and say how you will use partial results.

\exercisesection{Programming Exercises}

\pe Reproduce the degree table and \cref{fig:degree-curve} for T10. Then repeat
with a different random split of the $50$ points into $20/15/15$ ten times, and
report how often degree $3$ is chosen.

\pe Implement $k$-fold cross-validation from scratch and reproduce the five-fold
table. Repeat with $k = 10$ and with leave-one-out, and compare the chosen
degrees and the spreads.

\pe Simulate the bias--variance decomposition for ridge regression on the
degree-$10$ polynomial, for $\lambda$ from $10^{-4}$ to $10^2$, and plot bias$^2$,
variance and their sum against $\lambda$.

\pe Fit a logistic regression and an overfitted model (for instance, many
polynomial features with no penalty) to a classification dataset of your choice.
Draw a reliability diagram for each, report the Brier score and log loss, and
recalibrate the overfitted model by fitting a logistic regression to its logits
on validation data.

\exercisesection{Self-Test}

\qz High training error and similar validation error indicate: (a)~high
variance; (b)~high bias; (c)~a good model; (d)~a bug in the test set.

\qz The validation set is used to: (a)~fit the parameters; (b)~choose
hyperparameters and models; (c)~report final performance; (d)~nothing.

\qz In $5$-fold cross-validation each example is used for validation: (a)~never;
(b)~once; (c)~four times; (d)~five times.

\qz The noise term in the bias--variance decomposition: (a)~shrinks with more
data; (b)~shrinks with a better model; (c)~cannot be reduced by any model;
(d)~equals the variance.

\qz Increasing $\lambda$ tends to: (a)~increase variance; (b)~increase bias;
(c)~reduce both; (d)~have no effect.

\qz For a high-bias model, collecting more data will: (a)~help a lot;
(b)~not help much; (c)~always hurt; (d)~reduce bias to zero.

\qz Two classifiers with the same accuracy: (a)~must have the same log loss;
(b)~can differ in calibration; (c)~have the same reliability table; (d)~are
interchangeable.

\qz For a large neural network with high training error, the first remedy to try
is: (a)~more data; (b)~a bigger network; (c)~a larger $\lambda$; (d)~fewer
features.

\furtherreading

\textcite{james2021} §5.1 is the primary reading: it develops the validation-set
approach, leave-one-out and $k$-fold cross-validation, with a careful discussion
of the bias and variance of the cross-validation estimate itself (§5.1.4).
§2.2.2 states the bias--variance decomposition and illustrates it with the
U-shaped curves of this chapter. \textcite{bishop2006} §1.3 discusses model
selection and cross-validation briefly, and §3.2 derives the bias--variance
decomposition in the frequentist setting and illustrates it with regularized
fits to many independent datasets --- the simulation of \cref{sec:bias-variance}
done by hand, with plots well worth studying.

\bigskip
\noindent\emph{Next:} \Cref{ch:ch11} uses these diagnostics inside the iterative
loop of model development: error analysis, adding data, transfer learning,
metrics for skewed classes, and the fairness checks every deployed model needs.
